\documentclass[12pt]{article}

\usepackage[utf8]{inputenc}  
\usepackage[portuguese]{babel}   

\usepackage[ruled, vlined, algonl, portuguese]{algorithm2e}

\newcommand{\linesnumbered}{\LinesNumbered}
\newcommand{\dontprintsemicolon}{\DontPrintSemicolon}
\newcommand{\incmargin}{\IncMargin}
\newcommand{\decmargin}{\DecMargin}

\def\le{\leqslant}\def\leq{\le}
\def\ge{\geqslant}\def\geq{\ge}
\newcommand{\floor}[1]{\left\lfloor{#1}\right\rfloor}
\newcommand{\ceil}[1]{\left\lceil{#1}\right\rceil}
\newcommand{\pare}[1]{\left({#1}\right)}
\newcommand{\set}[1]{\left\{{#1}\right\}}
\newcommand{\range}[2]{\set{{#1},\dots,{#2}}}
\newcommand{\dist}{\mathrm{dist}}
\newcommand{\cluster}{\mathrm{cluster}}
\newcommand{\merge}{\mathrm{merge}}

\def\code#1{\texttt{#1}}


\title{Lista 01 de Análise de Algoritmos}
\date{}
\author{Turma do 4º ano}


\begin{document} 

\maketitle

\vspace{3em}


Para cada problema abaixo, desenvolva um algoritmo, 
descreva o algoritmo (pode ser em alto nível, pseudo-código ou 
linguagem de programação) e 
diga qual é a complexidade do seu algoritmo na notação $O$.

Utilize estratégia de programação dinâmica com memorização.

OBS: O algoritmo pode usar operações como saber o tamanho do 
vetor e pegar um elemento do vetor da posição $i$ sem custo de 
cálculo.

\begin{description}



\item[Fibonacci:] 
Os números de Fibonacci correspondem a uma sequência 
infinita na qual os dois primeiros termos são 0 e 1. 
Cada termo da sequência, 
à exceção dos dois primeiros, 
é igual à soma dos dois anteriores, 
conforme a relação de recorrência abaixo:
\[f_n = f_{n-1} + f_{n-2}\]
Desenvolva dois algoritmos que retorna o $i$-ésimo termo da sequência.



\item[Tribonacci:]
Seguindo o exercício anterior, 
escreva um algoritmo para o Tribonacci. 
Uma variação do Fibonacci em que o valor de um termo 
depende dos \textbf{3} termos anteriores.

\[f_n = f_{n-1} + f_{n-2} + f_{n-3},\]

e os três primeiros termos são 0, 0 e 1. 



\item[Subida das Escadas:]
Dada uma escada com $n$ degraus, uma pessoa parada em um degrau pode 
subir um ou dois degraus.
Qual a quantidade de combinações diferentes que uma pessoa pode usar para
subir a escada.



\item[Máximo Número de Segmentos:]
Dado um segmento de tamanho $n$. 
O objetivo é cortar o segmento na maior quantidade de partes possível.
Sendo que cada corte pode somente ter 3 tamanhos possíveis: $x$, $y$ e $z$.



\item[Roubo de Casas:] 
Dado um array de tamanho $n$, com números inteiros positivos.
Escolher posições no array tal que a soma dos seus elementos seja máxima, 
atendendo à restrição de que 
não se pode escolher duas posições vizinhas.
Caso a posição $i$ esteja entre as escolhidas, 
as posições $i-1$ e $i+1$ não poderão ser escolhidas.




\item[Troco em Moedas:]
Dada uma lista de tipos de \code{moedas[i]} de tamanho $n$ e 
um valor soma objetivo $S$, 
onde cada elemento \code{moedas[i]} é um valor diferente de moeda.
Considere que você tem uma quantidade infinita de cada moeda.
O objetivo é encontrar a mínima quantidade de moedas para juntar a soma $S$.
Se não for possível juntar a soma, retorne o número -1.




\item [Problema da Mochila:]
Dados dois arrays, \code{valor[]} e \code{peso[]}, 
onde cada elemento na posição \code{i} de \code{valor} representa o valor do item \code{i}, 
e cada elemento na posição \code{i} de \code{peso} representa o peso do item \code{i}.
Também é dado um valor \code{W} representando o peso total que a mochila aguenta.
O objetivo é juntar um conjunto de itens,
 tal que a soma tenha valor máximo e 
não ultrapasse o peso \code{W} da mochila.

OBS: Se um item for incluído, ele deve ser incluído por inteiro, não pode ser 
usado uma fração do item.



\end{description}


\end{document}

